\begin{thebibliography}{9}

\bibitem{homomorphic_encryption}
Gentry, C. (2009). Fully homomorphic encryption using ideal lattices. In \textit{Proceedings of the 41st annual ACM symposium on Theory of computing} (pp. 169-178). ACM.

\bibitem{tfhe_scheme}
Chillotti, I., Gama, N., Georgieva, M., \& Izabach\`ene, M. (2020). TFHE: Fast Fully Homomorphic Encryption over the Torus. \textit{Journal of Cryptology}, 33(1), 34-91.

\bibitem{tfhe_rs}
Zama. \textit{tfhe-rs: A Pure Rust Implementation of the TFHE Scheme}. Retrieved from \url{https://github.com/zama-ai/tfhe-rs}. All evaluation measurements in this paper were taken against version 1.5.0 with default parameters.

\bibitem{zama_fhe}
Zama. (2024). Concrete: Fully Homomorphic Encryption (FHE) library in Rust. \textit{Zama}. Retrieved from \url{https://www.zama.ai/}

\bibitem{deterministic_tfhe}
Smart, N. P., \& Walter, M. (2025). Reactive Correctness, sINDCPA-D-Security and Deterministic Evaluation for TFHE. \textit{IACR Cryptology ePrint Archive}, Report 2025/2005. Retrieved from \url{https://eprint.iacr.org/2025/2005}. Shows that sINDCPA-D security ordinarily calls for a randomised evaluation procedure, and that the randomisation can be removed for TFHE in the random oracle model --- the security counterpart to the practical reproducibility requirement discussed in Section 3.2.

\bibitem{byzantine_generals}
Lamport, L., Shostak, R., \& Pease, M. (1982). The Byzantine Generals Problem. \textit{ACM Transactions on Programming Languages and Systems}, 4(3), 382-401.

\bibitem{shamir_secret_sharing}
Shamir, A. (1979). How to share a secret. \textit{Communications of the ACM}, 22(11), 612-613.

\bibitem{floating_point}
Goldberg, D. (1991). What Every Computer Scientist Should Know About Floating-Point Arithmetic. \textit{ACM Computing Surveys}, 23(1), 5-48.

\bibitem{cooley_tukey}
Cooley, J. W., \& Tukey, J. W. (1965). An algorithm for the machine calculation of complex Fourier series. \textit{Mathematics of Computation}, 19(90), 297-301.

\bibitem{wasm_spec}
World Wide Web Consortium. \textit{WebAssembly Core Specification}. Retrieved from \url{https://webassembly.github.io/spec/core/}

\bibitem{phala_network}
Phala Network. (2024). Phala Network: Privacy-preserving computation for Web3. \textit{Phala Network}. Retrieved from \url{https://phala.network/}. Phala provides cryptographic attestations that allow verification of the integrity of AI workloads in real-time, ensuring they run in genuine secure environments.

\end{thebibliography}
